% 第九卷 15.7.4 节：Moshi 在一个 80 ms 时间步内的 RQ-Transformer 计算（通栏，约 145 mm）。
% 蓝色：用户侧输入；黄色描边：两个 Transformer；红色：系统生成的词元与声音。
% 左侧时间 Transformer 每步运行一次，右侧深度 Transformer 在步内依次预测 K=17 个子词元。
\begin{tikzpicture}[
  x=1mm,y=1mm,
  figure picture,
  model/.style={figure node,draw=FigureModel,fill=FigureModelFill,
    minimum height=12mm,font=\figurefont\scriptsize},
  inp/.style={figure node,draw=FigureInput,fill=FigureInputFill,
    minimum height=8mm,font=\figurefont\scriptsize},
  tok/.style={figure node,draw=FigureOutput,fill=FigureOutputFill,
    minimum height=7mm,minimum width=12mm,inner xsep=1.5pt,
    font=\figurefont\scriptsize},
  utok/.style={tok,draw=FigureInput,fill=FigureInputFill},
  lab/.style={font=\figurefont\scriptsize,text=FigureMuted,align=center},
  fl/.style={figure flow},
  back/.style={figure flow,draw=FigureMuted,dash pattern=on 2pt off 1.5pt}
]
  % 时间维：每 80 ms 运行一次的大模型。
  \node[inp,minimum width=36mm] (hist) at (20,0) {历史步$V_0,\ldots,V_{s-1}$\\每步17个词元};
  \node[model,minimum width=36mm] (emb) at (20,15) {每步17个嵌入相加\\得到$\mathbb R^{4096}$输入};
  \node[model,minimum width=36mm,minimum height=16mm] (temp) at (20,34)
    {时间Transformer\\Helium初始化\\32层，$d=4096$};
  \draw[fl] (hist) -- (emb);
  \draw[fl] (emb) -- (temp);

  % 深度维：步内依次预测 17 个子词元的小模型。
  \node[model,minimum width=84mm,minimum height=14mm] (depth) at (100,34)
    {深度Transformer（6层，$d=1024$，各$k$有独立投影）\\读取$\bm z_s$与本步已生成的$V_{s,1},\ldots,V_{s,k-1}$，输出$V_{s,k}$的分布};
  \draw[fl] (temp.east) -- node[figure edge label,above=0.5mm] {$\bm z_s$} (depth.west);

  % 第 s 步的 17 个子词元。
  \node[tok] (w) at (64,58) {$W_s$\\文字};
  \node[tok] (a1) at (79,58) {$A_{s,1}$\\语义};
  \node[tok,minimum width=19mm] (aq) at (97,58) {$A_{s-\tau,2:8}$\\声学7级};
  \node[utok] (u1) at (116,58) {$A'_{s,1}$\\用户语义};
  \node[utok,minimum width=19mm] (uq) at (134,58) {$A'_{s-\tau,2:8}$\\用户声学};
  \draw[fl] (w) -- (a1);
  \draw[fl] (a1) -- (aq);
  \draw[fl] (aq) -- (u1);
  \draw[fl] (u1) -- (uq);
  \foreach \n in {w,a1,aq,u1,uq} {\draw[fl] (depth.north -| \n) -- (\n.south);}
  \node[figure edge label,font=\figurefont\scriptsize] at (64,47) {$k=1$};
  \node[figure edge label,font=\figurefont\scriptsize] at (79,47) {$k=2$};
  \node[figure edge label,font=\figurefont\scriptsize] at (97,47) {$k=3,\ldots,9$};
  \node[figure edge label,font=\figurefont\scriptsize] at (125,47) {$k=10,\ldots,17$};

  % 输出与输入的声音。
  \node[figure node,draw=FigureOutput,fill=FigureOutputFill,minimum width=32mm,
    font=\figurefont\scriptsize] (dec) at (88,82) {Mimi解码器\\系统声音，80\,ms一帧};
  \draw[fl] (a1.north) -- (dec.south -| a1);
  \draw[fl] (aq.north) -- (dec.south -| aq);
  \node[inp,minimum width=36mm] (enc) at (125,82) {用户声音$\to$Mimi编码器\\推断时取实际输入};
  \draw[fl,draw=FigureInput] (enc.south -| u1) -- (u1.north);
  \draw[fl,draw=FigureInput] (enc.south -| uq) -- (uq.north);

  % 本步结果写回历史。
  \draw[back] (w.west) -- (50,58) |- (-4,72) -- (-4,0) -- (hist.west);
  \node[lab,anchor=south west] at (-3,72.5) {第$s$步的17个词元写回历史，供下一步读取};
\end{tikzpicture}
